% Generated draft: requires mathematical and source review.
\documentclass[11pt]{article}
\usepackage[margin=25mm]{geometry}
\usepackage{fontspec}
\setmainfont{Latin Modern Roman}
\usepackage{amsmath,amssymb,mathtools,unicode-math}
\setmathfont{Latin Modern Math}
\usepackage{enumitem,needspace,xurl,hyperref}
\setlength{\parindent}{0pt}
\setlength{\parskip}{5pt}
\setlength{\emergencystretch}{3em}
\newcommand{\N}{\mathbb{N}}
\newcommand{\Z}{\mathbb{Z}}
\newcommand{\Q}{\mathbb{Q}}
\newcommand{\R}{\mathbb{R}}
\newcommand{\C}{\mathbb{C}}
\newcommand{\F}{\mathbb{F}}
\newcommand{\A}{\mathbb{A}}
\newcommand{\PP}{\mathbb{P}}
\newcommand{\E}{\mathbb{E}}
\newcommand{\eps}{\varepsilon}
\newcommand{\dd}{\,\mathrm{d}}
\begin{document}
\section*{UnsolvedMath ID: 30002285}
\textbf{Source problem code:} catalogue entry
\subsection{Definitions and mathematical statement}
Let \(G\) be a fixed finite group, and let \(\rho\) denote the topological action type appearing in the supplied source record. Write \(\mathcal M_{g,G,\rho}\) for the source's moduli space of curves of genus \(g\) with a \(G\)-action of type \(\rho\). For a represented curve \(C\), let \(g'\) denote the genus of the quotient curve \(C/G\).

The supplied conjectural question is:
\begin{quote}
For fixed \(G\) and action data \(\rho\), do the cohomology groups of \(\mathcal M_{g,G,\rho}\) stabilize as the quotient genus \(g'\) tends to infinity?
\end{quote}
Here \(g'\), rather than the genus \(g\) of the covering curve, is the stated stabilization parameter.

This question is not a fully specified mathematical proposition in the supplied record. In particular, a topological action type on a fixed surface does not itself specify action types at arbitrarily large quotient genus. A family of action types extending \(\rho\), and the branch or other action data held fixed in that family, must therefore be specified. The record also does not specify the class of curves, the base field, whether \(\mathcal M_{g,G,\rho}\) is a coarse moduli space or a moduli stack, the cohomology theory and coefficients, or comparison maps.

For clarity, one possible interpretation of the word ``stabilize'' would be eventual constancy in each fixed cohomological degree: after specifying a family \(X_{g'}\) of these moduli objects, a cohomology theory, and a coefficient object \(A\), one would ask whether
\[
\forall k\in\Z_{\geq 0}\ \exists N_k\in\Z_{\geq 0}\ \forall g'_1,g'_2\geq N_k,
\qquad H^k(X_{g'_1};A)\cong H^k(X_{g'_2};A).
\]
This is only an illustrative interpretation, not an additional assertion attributed to the source. The record does not determine whether stabilization means such abstract isomorphisms, isomorphisms induced by specified comparison maps, or eventual stabilization of the entire graded cohomology. No choice among these interpretations is made here.
\subsection{Short English statement}
The supplied question asks whether the cohomology of moduli spaces of curves with a fixed finite-group action stabilizes when the genus of the quotient curve grows. A precise formulation requires a compatible family of action types, a specified moduli object, cohomology theory and coefficients, and a definition of stabilization. The supplied record labels the problem open; this draft does not independently verify that status.
\subsection{Sources}
\begin{itemize}
\item Local UnsolvedMath registry: supplied complete record for permanent ID \(30002285\), titled ``Cohomological Stability for Moduli of Curves with Group Actions,'' in the Oberwolfach Reports Open Problems collection; exact source problem code \texttt{catalogue entry}. The record gives the proposed year as 2013 and names no individual proposer.
\item Supplied report citation: \emph{Moduli Spaces in Algebraic Geometry} (2013), Oberwolfach Reports, pp.\ 343--392, DOI \texttt{10.4171/owr/2013/06}. \url{https://doi.org/10.4171/owr/2013/06}. The supplied background describes Oberwolfach Reports as documenting mathematical workshops and their open-problem discussions. No specific page or passage for this question is supplied.
\item Supplied literature reference: Catanese--Loenne--Perroni, \emph{Genus stabilization}. \url{https://arxiv.org/abs/2508.10199}.
\item Literature assessment dated exactly \texttt{2026-08-21}: the supplied curation labels the problem open, states that component stabilization is known, and reports homology stabilization only for a framed variant with at most one branch point, while describing the general cohomology-stabilization conjecture as unresolved. This is supplied dated curation, not an independently verified literature assessment. The record itself requests independent verification of time-sensitive status and cited claims.
\item The registry labels its statement \texttt{corrected\_verified} and reports that the official report identifies the moduli objects as curves with a fixed finite-group action type and the stabilization parameter as quotient-curve genus. It says this context was restored and the extracted conjecture label removed. These verification claims are supplied metadata; the original extracted statement nevertheless begins with ``Conjecture.'' Neither source inspection nor independent verification is asserted here.
\end{itemize}
\end{document}
